\documentclass[12pt]{article}
\usepackage[a4paper,margin=1in]{geometry}
\usepackage{amsmath,amssymb}
\usepackage{mathtools}
\usepackage{bm}

\title{Monte Carlo Wave-Function Method in QuTiP for the Interacting Dissipative Single-Mode Bosonic System}
\author{}
\date{}

\begin{document}
\maketitle

\section*{Purpose}
This note gives a thesis-friendly explanation of what the Monte Carlo wave-function method does mathematically in QuTiP for the present single-mode bosonic problem with nonzero interaction strength $U$ and nonzero loss rate $\gamma$. The goal is to connect the model Hamiltonian, the Lindblad equation, the stochastic trajectory picture, and the actual code used in this project.

\section*{Model}
We consider a single bosonic mode with annihilation operator $\hat a$ and number operator
\[
\hat n = \hat a^\dagger \hat a.
\]
The interacting Hamiltonian is
\[
\hat H = \frac{U}{2}\hat a^\dagger \hat a^\dagger \hat a \hat a
= \frac{U}{2}\hat n(\hat n-1),
\]
and the single-particle loss channel is represented by the collapse operator
\[
\hat C = \sqrt{\gamma}\,\hat a.
\]

The density operator $\rho$ obeys the Lindblad master equation
\[
\frac{d\rho}{dt}
= -i[\hat H,\rho]
+ \hat C \rho \hat C^\dagger
- \frac12\{\hat C^\dagger \hat C,\rho\}.
\]
Since
\[
\hat C^\dagger \hat C = \gamma \hat a^\dagger \hat a = \gamma \hat n,
\]
this becomes
\[
\frac{d\rho}{dt}
= -i[\hat H,\rho]
+ \gamma \hat a \rho \hat a^\dagger
- \frac{\gamma}{2}\{\hat n,\rho\}.
\]

\section*{Why Monte Carlo Trajectories Can Replace the Master Equation}
Solving the master equation evolves the full density matrix, whose size grows rapidly with Hilbert-space truncation. The Monte Carlo wave-function method rewrites the same open-system dynamics as an ensemble average over stochastic pure-state trajectories. Instead of evolving $\rho(t)$ directly, one evolves many state vectors
\[
|\psi_r(t)\rangle,
\qquad r=1,\dots,N_{\mathrm{traj}},
\]
and approximates expectation values by trajectory averaging:
\[
\langle \hat O(t)\rangle
= \mathrm{Tr}\!\bigl(\rho(t)\hat O\bigr)
\approx \frac{1}{N_{\mathrm{traj}}}\sum_{r=1}^{N_{\mathrm{traj}}}
\langle \psi_r(t)|\hat O|\psi_r(t)\rangle.
\]

Mathematically, the method is exact in the limit of infinitely many trajectories:
\[
\rho(t) = \mathbb{E}\bigl[|\psi(t)\rangle\langle\psi(t)|\bigr].
\]
Thus the Monte Carlo solver does not change the physics; it changes the numerical representation of the same Lindblad evolution.

\section*{Effective Non-Hermitian Hamiltonian}
Between jumps, each trajectory evolves under the effective non-Hermitian Hamiltonian
\[
\hat H_{\mathrm{eff}}
= \hat H - \frac{i}{2}\hat C^\dagger \hat C.
\]
For the present model,
\[
\hat H_{\mathrm{eff}}
= \frac{U}{2}\hat a^\dagger \hat a^\dagger \hat a \hat a
- \frac{i\gamma}{2}\hat a^\dagger \hat a
= \frac{U}{2}\hat n(\hat n-1) - \frac{i\gamma}{2}\hat n.
\]

The Schr\"odinger-type equation for a single trajectory between jumps is
\[
\frac{d}{dt}|\psi(t)\rangle
= -i\hat H_{\mathrm{eff}}|\psi(t)\rangle.
\]
Because $\hat H_{\mathrm{eff}}$ is non-Hermitian, the norm of the unnormalized state decreases:
\[
\frac{d}{dt}\langle \psi(t)|\psi(t)\rangle \le 0.
\]
This norm loss has a direct probabilistic meaning: it encodes the probability that no quantum jump has yet occurred.

\section*{Physical Meaning of the Two Possible Updates}
At the trajectory level, the dynamics consists of two alternating processes.

\subsection*{1. No-jump evolution}
If no photon-loss event occurs during a short interval $dt$, the state evolves according to
\[
|\tilde\psi(t+dt)\rangle
\approx \left(1 - i\hat H_{\mathrm{eff}}dt\right)|\psi(t)\rangle.
\]
The state is then renormalized:
\[
|\psi(t+dt)\rangle
= \frac{|\tilde\psi(t+dt)\rangle}{\sqrt{\langle \tilde\psi(t+dt)|\tilde\psi(t+dt)\rangle}}.
\]

For small $dt$, the no-jump probability is
\[
p_{\mathrm{no\ jump}}(t,t+dt)
= 1 - dt\,\langle \psi(t)|\hat C^\dagger \hat C|\psi(t)\rangle + O(dt^2).
\]
Since $\hat C^\dagger \hat C=\gamma \hat n$, this becomes
\[
p_{\mathrm{no\ jump}}(t,t+dt)
= 1 - \gamma\,dt\,\langle \hat n\rangle_t + O(dt^2).
\]
So the probability of a jump is larger when the mode occupation is larger.

\subsection*{2. Jump evolution}
If a loss event occurs, the jump operator is applied:
\[
|\psi(t)\rangle \longrightarrow
\frac{\hat C|\psi(t)\rangle}{\sqrt{\langle \psi(t)|\hat C^\dagger \hat C|\psi(t)\rangle}}
= \frac{\hat a|\psi(t)\rangle}{\sqrt{\langle \hat n\rangle_t}}.
\]
For a single collapse channel there is no ambiguity about which jump occurs: every jump corresponds to one action of $\hat a$, meaning the loss of one boson from the mode.

\section*{How the Randomness Enters}
The stochastic part of the method comes from random numbers used to decide when jumps occur. Conceptually, one may describe the algorithm over a small time step $dt$ as follows:
\begin{enumerate}
\item start from a normalized state $|\psi(t)\rangle$,
\item compute the jump probability
\[
p_{\mathrm{jump}} \approx dt\,\langle \psi(t)|\hat C^\dagger \hat C|\psi(t)\rangle
= \gamma\,dt\,\langle \hat n\rangle_t,
\]
\item draw a random number $r\in(0,1)$,
\item if $r<p_{\mathrm{jump}}$, perform a jump,
\item otherwise evolve with $\hat H_{\mathrm{eff}}$ and renormalize.
\end{enumerate}

This is the intuitive discrete-time picture and it is useful for explaining the method in words.

\section*{What QuTiP Actually Does Internally}
QuTiP's \texttt{mcsolve} implements the same physics more accurately than the simple fixed-step picture above. Instead of deciding jumps only at the ends of user-defined output steps, it propagates the state continuously under $\hat H_{\mathrm{eff}}$ and monitors the norm of the unnormalized state.

For each trajectory, QuTiP samples a random threshold
\[
r_1 \in (0,1),
\]
then evolves the unnormalized state until
\[
\|\tilde\psi(t)\|^2 = r_1.
\]
The time at which this equality is reached is the collapse time. QuTiP therefore performs an internal root-finding step to locate the jump time accurately. Once the collapse time is found, the jump operator is applied and the state is renormalized.

Because there is only one collapse operator in this project,
\[
\hat C = \sqrt{\gamma}\,\hat a,
\]
every collapse is automatically of that type. In a multi-channel problem QuTiP would also choose which collapse operator acts according to relative jump probabilities, but that extra random choice is not needed here.

This detail is important for thesis writing: the user-defined spacing $dt$ in the code sets the list of times at which observables are recorded, but QuTiP may use smaller internal integration steps and a separate collapse-time search between those output times.

\section*{Role of Nonzero $U$ and Nonzero $\gamma$}
When both $U\neq 0$ and $\gamma\neq 0$, the trajectory evolution contains two competing effects:
\begin{itemize}
\item the interaction term
\[
\frac{U}{2}\hat a^\dagger \hat a^\dagger \hat a \hat a
\]
produces number-dependent phase evolution,
\item the dissipative term
\[
-\frac{i\gamma}{2}\hat a^\dagger \hat a
\]
causes norm decay in the no-jump branch and therefore generates stochastic jumps.
\end{itemize}

So the interacting Monte Carlo trajectory is not just exponential decay. Between jumps, each number component $|n\rangle$ acquires both
\begin{itemize}
\item a phase factor from the Kerr interaction,
\item and an amplitude damping factor from the imaginary part of $\hat H_{\mathrm{eff}}$.
\end{itemize}
Then, at random collapse times, the action of $\hat a$ moves weight from $|n\rangle$ to $|n-1\rangle$. Repeating this many times produces a stochastic history of losses superimposed on coherent interacting evolution.

\section*{Trajectory Averaging}
One trajectory is only one possible realization of the open-system dynamics. The physical observables are obtained by averaging many such realizations. For an operator $\hat O$,
\[
\langle \hat O(t)\rangle
\approx \frac{1}{N_{\mathrm{traj}}}\sum_{r=1}^{N_{\mathrm{traj}}}
\langle \psi_r(t)|\hat O|\psi_r(t)\rangle.
\]
As $N_{\mathrm{traj}}$ increases, the statistical noise decreases roughly like
\[
\frac{1}{\sqrt{N_{\mathrm{traj}}}}.
\]
This is why Monte Carlo methods trade density-matrix size for sampling noise: one avoids evolving the full density matrix, but one must average many trajectories to obtain smooth results.

\section*{How It Was Coded in This Project}
The implementation is in the file \texttt{src/bosonic\_dissipation/monte\_carlo\_method.py}. The coded procedure is:

\subsection*{1. Build the Hilbert space and initial state}
The annihilation operator is created on a truncated Fock space:
\[
\hat a = \texttt{qt.destroy(hilbert\_size)}.
\]
Then the initial state is chosen as either:
\begin{itemize}
\item a Fock state $\lvert n_0\rangle$, or
\item a coherent state $\lvert \alpha_0\rangle$ with $\alpha_0=\sqrt{n_0}$.
\end{itemize}

\subsection*{2. Build the operators}
The code constructs
\[
\hat n = \hat a^\dagger \hat a,
\qquad
\hat F_2 = \hat a^{\dagger 2}\hat a^2,
\]
and the Hamiltonian
\[
\hat H = \frac{U}{2}\hat a^{\dagger 2}\hat a^2.
\]
The collapse-operator list is
\[
\texttt{collapse\_operators = [}\sqrt{\gamma}\,\hat a\texttt{]}.
\]

\subsection*{3. Build the output-time grid}
The code defines
\[
t_0, t_1, \dots, t_M
\]
using
\[
\texttt{time\_values = np.arange(0.0, time + dt, dt)}.
\]
These are the times at which QuTiP returns expectation values. They should be interpreted as output sampling times, not necessarily as QuTiP's internal solver step size.

\subsection*{4. Run QuTiP's Monte Carlo solver}
The core call is
\begin{verbatim}
result = qt.mcsolve(
    hamiltonian,
    psi0,
    time_values,
    collapse_operators,
    [a, n_op, factorial_second_moment_op],
    ntraj=num_of_samples,
    seeds=seed,
    options={"progress_bar": False},
)
\end{verbatim}
This tells QuTiP to:
\begin{itemize}
\item evolve \texttt{num\_of\_samples} stochastic trajectories,
\item use the interacting Hamiltonian $\hat H$,
\item include the loss process through $\sqrt{\gamma}\,\hat a$,
\item and record, at each requested time, the expectation values of
\[
\hat a, \qquad \hat n, \qquad \hat a^{\dagger 2}\hat a^2.
\]
\end{itemize}

\subsection*{5. Average over trajectories and post-process observables}
After \texttt{mcsolve} finishes, QuTiP has already averaged the requested expectation values over all trajectories. The code then interprets them as:
\[
g^{(1)}(t) \equiv \langle \hat a(t)\rangle,
\]
\[
\langle \hat n(t)\rangle,
\]
\[
F_2(t) \equiv \langle \hat a^{\dagger 2}\hat a^2\rangle
= \langle \hat n(\hat n-1)\rangle.
\]

From these stored arrays, the code computes the variance through
\[
\mathrm{Var}(\hat n)
= \langle \hat n^2\rangle - \langle \hat n\rangle^2
= F_2 + \langle \hat n\rangle - \langle \hat n\rangle^2,
\]
because
\[
\hat a^{\dagger 2}\hat a^2 = \hat n(\hat n-1).
\]
It also computes
\[
g^{(2)}(t)
= \frac{\langle \hat a^{\dagger 2}\hat a^2\rangle}{\langle \hat n\rangle^2}
= \frac{F_2(t)}{\langle \hat n(t)\rangle^2},
\]
whenever the denominator is nonzero.

\section*{A Thesis-Friendly Step-by-Step Summary}
The full method can be summarized in words as follows:
\begin{enumerate}
\item Start from the Hamiltonian
\[
\hat H = \frac{U}{2}\hat a^{\dagger 2}\hat a^2
\]
and the collapse operator
\[
\hat C=\sqrt{\gamma}\,\hat a.
\]
\item Replace the Lindblad master equation by a stochastic pure-state evolution.
\item For one trajectory, propagate the state under
\[
\hat H_{\mathrm{eff}}=\hat H-\frac{i}{2}\hat C^\dagger \hat C
\]
until a randomly determined jump time is reached.
\item At that random time, apply the jump operator $\hat a$ and renormalize the state.
\item Continue this no-jump/jump alternation throughout the full simulation interval.
\item Repeat the whole process for many independent trajectories.
\item Average the trajectory expectation values to approximate the physical observables of the open quantum system.
\end{enumerate}

\section*{Interpretation}
Mathematically, the Monte Carlo method used here is therefore an \emph{unraveling} of the Lindblad equation into random quantum-jump histories. The interaction strength $U$ enters through the deterministic but non-Hermitian inter-jump evolution, while the loss rate $\gamma$ determines the jump statistics. The final observables are not taken from one trajectory but from an ensemble average over many trajectories, which is exactly how QuTiP's \texttt{mcsolve} is used in the project.

\section*{Conclusion}
For the present single-mode interacting dissipative bosonic system, QuTiP's Monte Carlo solver does the following: it evolves a pure state under an effective Hamiltonian, inserts random loss events according to the decay of the wavefunction norm, applies the collapse operator $\sqrt{\gamma}\,\hat a$ at the collapse time, and then averages the resulting observables over many trajectories. In this way, it reproduces the same Lindblad dynamics that would otherwise be obtained from the density-matrix master equation, but through stochastic trajectory sampling.

\end{document}
